\chapter{Скільки потрібно входів}
% full version: chapters/19_dendrites.tex

Нейрон має розгалужені вхідні відростки --- дендрити. В одних нейронів це ціле дерево з тисячами контактів, в інших --- одна-дві гілки або й нічого. Для інженера кількість входів --- параметр схеми, і його підібрано під задачу.

\section*{Від нуля до ста тисяч}

У датчиків дотику й болю дендритів немає зовсім: провід просто тягнеться від шкіри до спинного мозку, а датчик сидить на самому його кінці. Це лінія з датчиком на дальньому краю.

\pencilfig{19}{\pencilart{19}}{Багато входів --- суматор; один величезний вхід --- швидкий і точний ретранслятор.}

У нейронів нюху й у нейронів, що передають сигнал від датчиків ока та вуха, --- один дендрит. Один вхід, один вихід: повторювач, буфер.

У детекторів напрямку звуку --- два дендрити, що дивляться в різні боки, по одному на кожне вухо. Рознісши два входи по різних гілках, нейрон робить «і» суворішим: дві порції струму в одній гілці заважали б одна одній.

У дрібних нейронів схеми навчання рухів --- близько чотирьох коротких дендритів, кожен хапає свій вхід. Таких нейронів десятки мільярдів, і кожен додає лише чотири сигнали. Зате разом вони розкладають вхід на мільйони комбінацій, а великий вихідний нейрон із сотнею тисяч входів вибирає з них потрібні.

І нарешті, нейрони з величезними деревами й тисячами входів --- суматори й класифікатори, у яких окремі гілки працюють майже як самостійні маленькі схеми.

\section*{Чому так}

Усе вирішує заряджання. Маленький нейрон із небагатьма входами заряджається від одного великого з'єднання за частки мілісекунди: такий нейрон точний у часі й надійно передає сигнал далі. Велике дерево одне з'єднання не зарядить ніколи --- заряд витече раніше. Йому потрібні десятки сигналів, що прийшли разом, зате воно може їх додавати й порівнювати. Мало входів і великі з'єднання --- для точного часу. Багато входів --- для лічби.

\pencilfig{128}{%
% charging: a small neuron reaches threshold from one large synapse at once; a big tree barely moves
% from one input and needs many inputs arriving together
\foreach \dx/\t in {0/{мало входів}, 4.3/{велике дерево}}{
  \draw[pen] (\dx,0) -- (\dx+3.6,0);
  \draw[pcGray, densely dashed, line width=0.5pt] (\dx,0.9) -- (\dx+3.6,0.9);
  \node[font=\small\itshape, text=pcGray, anchor=south] at (\dx+1.8,1.75) {\t};}
\node[lbl, anchor=east] at (-0.05,0.9) {поріг};
% small neuron
\draw[pcBlue, line width=2pt] (0.6,-0.15) -- (0.6,-0.6);
\draw[penlite=pcBlue] (0,0) -- (0.6,0) -- plot[domain=0.6:0.8, samples=12] (\x,{1.3*(1-exp(-(\x-0.6)/0.18))}) -- (0.8,1.6) -- (0.84,0) -- (3.5,0);
\node[lbl, text=pcRed, anchor=west] at (0.85,1.45) {клацання одразу};
\node[lbl, text=pcBlue, anchor=north] at (0.6,-0.62) {один великий вхід};
% big tree
\draw[pcBlue, line width=0.6pt] (4.8,-0.15) -- (4.8,-0.45);
\foreach \s in {6.15,6.2,6.25,6.3,6.35,6.4,6.45,6.5}{\draw[pcBlue, line width=0.6pt] (\s,-0.15) -- (\s,-0.45);}
\draw[penlite=pcBlue] (4.3,0) -- (4.8,0) -- plot[domain=4.8:6.15, samples=20] (\x,{0.16*((\x-4.8)/0.15)*exp(1-(\x-4.8)/0.15)}) -- plot[domain=6.15:6.62, samples=14] (\x,{1.25*(1-exp(-(\x-6.15)/0.4))}) -- (6.62,1.6) -- (6.66,0) -- (7.9,0);
\node[lbl, anchor=south] at (4.95,0.2) {майже нічого};
\node[lbl, text=pcBlue, anchor=north] at (6.33,-0.47) {багато разом};
}{Маленький нейрон заряджається від одного великого з'єднання за частки мілісекунди й одразу клацає. Велике дерево від одного входу майже не змінюється: йому потрібні десятки сигналів, що прийшли разом.}

Так ми отримали третє сортування нейронів: за речовиною, за приладом і за кількістю входів. Кожне показує ту саму машину з нового боку.

\fullversion{19}
